% ---
% id: receiver/rem:unchecked-assumptions
% kind: rem
% title: Assumptions this module cannot check
% summary: >-
%   The assumptions on the physical receiver and the run, and [[A-EIG]] on the
%   arithmetic, are never verified. The conditions decided on the declared data
%   are definitions with requirements, and the floating-point evaluation is a
%   restriction the contract places on itself.
% prominence: standard
% status: accepted
% uses: []
% mentions: [ass:eig, def:graph, req:graph-refusal, def:probe, req:probe-diagonal-admission, def:evaluation-model, req:eig-ledger-row]
% assumes: []
% refs: []
% provenance: newly-derived
% lean:
% ---

\begin{remark}\label{rem:unchecked-assumptions}
Every assumption of this contract that is a statement about the physical
receiver or the run is one the engine never verifies, as the contract says
of each: no function of the \tref{declared data} decides it, and validation
decides only those of its clauses that constrain the declared data.
Assumption~\ref{ass:eig} is a statement about the \emph{arithmetic}; it is
likewise never verified, and it is invoked only at a setting where a caller
declares an \tref{exact-window depth} and the weight is evaluated in floating
point (Requirement~\ref{req:eig-ledger-row}), so a declaration that declares
none, or whose weights are all evaluated rigorously, is unaffected by it. The conditions that constrain the declared data alone are not
assumptions of this contract: the well-formedness of a declared element graph
(Definition~\ref{def:graph}) and the occupation-diagonality of the declared
probe states (Definition~\ref{def:probe}) are decided on the declared data
(Requirements~\ref{req:graph-refusal} and~\ref{req:probe-diagonal-admission}).
The floating-point evaluation of Definition~\ref{def:evaluation-model} is a
restriction this contract places on itself and asserts nothing that a run
could violate.
\end{remark}
